% Part: applied-modal-logics
% Chapter: temporal-logic
% Section: introduction

\documentclass[../../../include/open-logic-section]{subfiles}

\begin{document}

\olfileid[id]{aml}{tl}{int}

\olsection{Pengantar}

Logika temporal membahas pernyataan tentang hal-hal yang akan atau pernah
berlaku. Arthur Prior dianggap sebagai penggagas logika temporal, yang ia
sebut \emph{logika kala (tense logic)}. Pembahasan logika temporal
kita di sini sebagian besar mengikuti pendekatan modal asli Prior: operator
temporal diperkenalkan ke dalam kerangka dasar logika proposisional, yang
memperlakukan pernyataan seolah-olah pada umumnya tidak berpenanda kala.

Sebagai contoh, dalam logika proposisional, saya mungkin berbicara tentang
seekor anjing bernama Beezie, yang kadang duduk dan kadang tidak duduk,
sebagaimana lazimnya anjing. Dalam logika klasik, menyatakan bahwa Beezie
sedang duduk sekaligus bahwa Beezie tidak sedang duduk merupakan
kontradiksi. Akan tetapi, jelas bahwa keduanya dapat benar, hanya saja tidak
pada waktu yang sama; penambahan operator temporal ke dalam bahasa
memungkinkan kita menyatakan hal itu dengan relatif mudah. Penambahan
operator temporal juga memungkinkan kita menjelaskan validitas inferensi
seperti inferensi dari ``Beezie akan mendapat camilan atau bola'' menuju
``Beezie akan mendapat camilan atau Beezie akan mendapat bola.''

Namun, logika temporal menimbulkan banyak persoalan filosofis yang dapat
mendorong kita memilih satu kerangka logika temporal alih-alih kerangka
lain. Sebagai contoh, suatu kontingen masa depan ialah pernyataan
tentang masa depan yang tidak niscaya dan juga tidak mustahil. Jika kita
mengatakan ``Richard akan pergi ke toko bahan makanan besok,'' kita
menyatakan sesuatu yang belum terjadi dan yang nilai kebenarannya dapat
diperdebatkan. Bahkan, dapat diperdebatkan apakah pernyataan itu sejak awal
dapat \emph{diberi} nilai kebenaran. Jika kita merupakan determinis ketat,
kita mungkin dapat menerima gagasan bahwa kalimat itu sebenarnya benar atau
salah bahkan sebelum peristiwa yang dimaksud seharusnya berlangsung---kita
mungkin hanya belum mengetahui nilai kebenarannya. Sebaliknya, kita mungkin
meyakini masa depan yang sungguh terbuka, tempat nilai kebenaran kontingen
masa depan belum ditentukan.

Ternyata, banyak komitmen mengenai struktur dan hakikat waktu ini tertanam
dalam pilihan model dan kerangka logika temporal kita. Sebagai contoh, kita
dapat bertanya apakah model yang kita bangun harus menggambarkan waktu
sebagai linear, bercabang, atau bahkan melingkar. Kita mungkin harus
memutuskan apakah model temporal kita mempunyai titik awal dan titik akhir,
serta apakah waktu direpresentasikan dengan saat-saat diskret atau sebagai
suatu kontinuum.

\end{document}
